\lecture{3}
\title{OS Design}

\begin{document}

\subsection{Design Ideals}

One of the main ideals of OS design is \textbf{isolation}.
Without isolation, issues like these might arise:
\begin{itemize}
    \item Suppose both the shell \verb|sh| and the cat program \verb|cat|
    have instructions that live in RAM.
    
    If \verb|sh| lives at address \verb|0x1000|,
    then perhaps the \verb|cat| program can pipe
    \verb|st 5, x1000| into the shell,
    thereby writing over the instructions that make the shell work.
    \item Or more simply,
    maybe one application uses too much memory,
    or hogs the CPU,
    or writes into another process' memory.
\end{itemize}

Some additional things we need to consider:
\begin{itemize}
    \item Every program---whether it's \verb|cat| or \verb|sh|---uses
    the same binary-file structure:
    \verb|.text| instructions at \verb|x1000|,
    followed by \verb|.data|, and so on.
    How can so many processes all have a \verb|x1000|?
    \item We need to implement some filesystem \verb|fs|
    that offers us convenience in referring to
    individual processes.
    Instead of saying, ``run the process at \verb|x2000| in RAM'',
    it would be nice to just say, ``run the process named \verb|cat|''.
\end{itemize}
In summary, the goal of the kernel designer
is to make sure that no user application can mess with the kernel.

\subsection{Hardware Support for Isolation}

Here, we outline the main ideas of relevant bits of hardware.
Future lectures will be dedicated to the implementation details.

\textbf{Page Tables.}
These translate virtual addresses into physical addresses.
This resolves the issue of giving every process its own \verb|x1000|---that is,
it gives every process its own memory.

\textbf{Modes.}
User applications (e.g., \verb|sh| or \verb|cat|) run in \textbf{user mode},
whereas the CPU and RAM run in \textbf{kernel mode}.
For example, modifying page tables should only be permitted in kernel mode.

\textbf{Transitioning.}
We need to be careful about transitioning between user mode and kernel mode.
When \verb|sh| or \verb|cat| call \verb|fork|,
it doesn't call \verb|fork| directly---instead, it calls
\verb|ECALL 1|, where the number indicates which system call is being executed.

When a user process calls \verb|ECALL|, it is \emph{not} a function calls---the
sort of calls where you move \verb|sp|, do some stuff, and then \verb|ra| back.
Instead, system calls are instructions to switch from user mode to kernel mode.
Similarly, to switch from kernel mode to user mode,
you run something called \verb|SRET|.

\begin{remark}
    Kernel mode is also sometimes called ``supervisor mode''.
    They're the same thing.

    We also user the terms ``user space'' and ``kernel space''
    to refer to processes / hardware that run in
    user mode or kernel mode.
    CPUs and RAM live in kernel space,
    whereas \verb|sh| and \verb|cat| live in user space.

    The OS 
\end{remark}

Interestingly, it is very hard to implement all of these perfectly.
The Linux kernel has roughly 40 million lines of code---and there are bugs.
People take this very seriously;
see \href{linuxcvetracker.com}{linuxcvetracker.com}.

\subsection{Kernel Mode}

A \textbf{monolothic OS} is an OS in which
all core OS services run in kernel mode
(including process scheduling, filesystems, e.g.).
This is in contrast to a \textbf{microkernel},
in which the bare minimum of OS services are in kernel mode
(e.g., address spaces thread scheduling),
whereas services like filesystems and network stacks
run as user-mode processes.

Some advantages:

\textbf{Performance.}
In a monolithic kernel,
a filesystem calling a disk driver is just a function call.
In a microkernel,
a filesystem calling a disk driver
involves 

\textbf{Code Length.}

\subsection{QEMU}

Note that \verb|qemu| is just a software implementation of
the RISC-V processor chip.
This board has 4 RISC-V cores (CPUs),
so if there are 8 processes,
then you might multiplex those 8 processes across 4 CPUs.

Even though we are interacting with software in this lab,
we should imagine that our software implementations
are just alterations to physical hardware on a RISC-V board.

\subsection{TODO}

The code for xv6 has most of its files organized into two directories:
\verb|user| and \verb|kernel|.
Since xv6 is a monolithic operating system,
all of the OS design is in \verb|kernel|.

\begin{remark}
    Allegedly, we are supposed to be able to
    read any line of the xv6 code
    and be able to explain exactly what it is doing.
\end{remark}

We might see, for example:

\begin{verbatim}
$ ls kernel
proc.c uart.c vm.c defs.h ...
\end{verbatim}

So each C file in the kernel implements a part of the OS
(e.g., \verb|proc.c| implements process management).
There are also header files;
for example, \verb|defs.h| has all of the APIs for all of the other C files
(e.g. processes, virtual memory, etc.).

If you run \verb|make kernel/kernel|,
this creates the assembly file \verb|kernel.asm|---that is,
you build the build.
The file \verb|kernel/kernel| is the compiled, linked xv6 kernel
that contains all \verb|.c| and \verb|.S| code in kernel,
combined into a single program.
It's a 64-bit RISC-V ELF executable;
it is the operating system.

RUnning \verb|make qemu| loads the file
into (emulated) RAM at address \verb|0x8000000|.
RISC-V will then run (emulated) instructions
starting at \verb|0x8000000|, and so on.

You can also run \verb|make kernel/kernel|
to compile each source file to a \verb|.o| file,
then links using \verb|kernel.ld|,
then generates \verb|kernel.asm| and \verb|kernel.sym|.
But this is not necessary;
running \verb|make qemu| recognizes the dependency on \verb|kernel/kernel|
and rebuilds things on its own.
Importantly, \verb|kernel.asm| and \verb|kernel.sym|
are not actually used; they're plain-text (ASCII) files
that exist solely for human debugging.

\subsection{TODO}

Note that \textbf{hart} is the exact same as \textbf{CPU} in RISC-V language.

Note that \verb|kernel/entry.S| contains \verb|.global _entry|,
which makes \verb|_entry| visible outside of \verb|entry.S|.
The only thing that \verb|entry.S| does is set up a stack.

The platform starts out as M mode (not supervisor mode),
which is the \emph{minimal} amount of code needed to get the platform going;
it says things like, if there's an interrupt, delegate to supervisor mode;
it has very little code, and so we never enter M mode (machine mode).
You can see this in \verb|start.c| in the kernel.

To get from M mode to supervisor mode, you run \verb|MRET|
(compare with \verb|SRET|).
So the end of \verb|start.c| is
\verb|asm volatile("mret");|.

Now let's discuss \verb|kernel/main.c|.
This sets up the device drivers (a bunch of \verb|init| files),
some disc IO, some virtual memory stuff, etc.
Notably, \verb|consoleinit();| is needed first;
when you run \verb|make qemu|, you see some things printed to output,
and this is because the line immediately after \verb|consoleinit();|
is \verb|printkinit();|.

Notably, only CPU 0 does the initialization code;
there's an \verb|if (cpuid() == 0)| at the very top
that ensures that initialization code is not run by multiple CPUs.

\begin{remark}
    It's called \verb|printk| because it's a kernel call.
    Using \verb|printf| creates a system call (\verb|ecall||)
    which traps into the kernel\dots but this only makes sense
    for calling \verb|printf| in user space;
    it would break if called in kernelspace.
\end{remark}

By looking at \verb|proc.h|, we can see what a process is.
It contains a bunch of registers,
a lock, a state, \dots not very much, actually,
only 15 or so lines of code in a structure.
The locks are there to coordinate how multiple CPUs
coordinate; we'll talk about locks next week.
There's also a \verb|pid| (the process ID),
a pagetable, open file descriptors,
the current working directory, and some other stuff.

In the kernel, files are \verb|inode|s,
because filesystem stuff.

In the HW, we asked:
\begin{itemize}
    \item What happens if we call fork
    but there's no procspace left because we're already running
    the maximum number of processes?
    \begin{itemize}
        \item It raises an error.
        Usually, fork calls \verb|allocproc|,
        which\dots we'll see.
    \end{itemize}
    \item If you call exec,
    it doesn't matter, because all exec does is
    replace the current process' memory with new memory.
    No new processes are needed, so it doesn't matter.
\end{itemize}

Every process in xv6 is created by fork except the very first one.
When fork is run, \verb|forkret| is run,
and this does some things,
including building a process memory with
\verb|.text| and \verb|.data| and so on,
and returns to userspace
by running \verb|kexec("/init", ...)|.

So what does \verb|user/init.c| do?
It sets up the environment for the shell to run.
It also \verb|dup(0)| twice to create stdin,
stdout, and stderr at descriptors 0, 1, and 2.
It then runs \verb|fork()|
to create a shell,
and prints some text that you see every time you run \verb|make qemu|.

Notably, register \verb|a7| is the one that
acts as a parameter for \verb|ECALL|.
We went from kernel mode to user mode
and from init we went back to kernel mode.

You can always modify the kernel too by putting printk statements
and seeing what happens.

With just one modification,
we can see how modifying \verb|syscall|,
that many, many system calls were necessary
to initiate the shell alone,
a lot of sys call 16s.
There's a sys call 1 from the fork.
The sys call 16s are from printing one character at a time
(it prints s, t, a, r, t, \dots and each character
is seen as a ssystem call).  Every character in a \verb|printf|
is a system call, and we print a lot of characters,
hence a lot of sys call 16s.

ts cooked

Main points:
\begin{itemize}
    \item KERNEL vs. USER space.
    \item System calls!  They're a little bit more concrete, hopefully.
    Useful in lab 2; we'll modify the syscall function, etc.
    \item The rest of the semester is just
    implementing the system calls,
    because there's not that many of them.
\end{itemize}

\end{document}